% Options for packages loaded elsewhere
\PassOptionsToPackage{unicode}{hyperref}
\PassOptionsToPackage{hyphens}{url}
\documentclass[
]{article}
\usepackage{xcolor}
\usepackage{amsmath,amssymb}
\setcounter{secnumdepth}{-\maxdimen} % remove section numbering
\usepackage{iftex}
\ifPDFTeX
  \usepackage[T1]{fontenc}
  \usepackage[utf8]{inputenc}
  \usepackage{textcomp} % provide euro and other symbols
\else % if luatex or xetex
  \usepackage{unicode-math} % this also loads fontspec
  \defaultfontfeatures{Scale=MatchLowercase}
  \defaultfontfeatures[\rmfamily]{Ligatures=TeX,Scale=1}
\fi
\usepackage{lmodern}
\ifPDFTeX\else
  % xetex/luatex font selection
\fi
% Use upquote if available, for straight quotes in verbatim environments
\IfFileExists{upquote.sty}{\usepackage{upquote}}{}
\IfFileExists{microtype.sty}{% use microtype if available
  \usepackage[]{microtype}
  \UseMicrotypeSet[protrusion]{basicmath} % disable protrusion for tt fonts
}{}
\makeatletter
\@ifundefined{KOMAClassName}{% if non-KOMA class
  \IfFileExists{parskip.sty}{%
    \usepackage{parskip}
  }{% else
    \setlength{\parindent}{0pt}
    \setlength{\parskip}{6pt plus 2pt minus 1pt}}
}{% if KOMA class
  \KOMAoptions{parskip=half}}
\makeatother
\usepackage{graphicx}
\makeatletter
\newsavebox\pandoc@box
\newcommand*\pandocbounded[1]{% scales image to fit in text height/width
  \sbox\pandoc@box{#1}%
  \Gscale@div\@tempa{\textheight}{\dimexpr\ht\pandoc@box+\dp\pandoc@box\relax}%
  \Gscale@div\@tempb{\linewidth}{\wd\pandoc@box}%
  \ifdim\@tempb\p@<\@tempa\p@\let\@tempa\@tempb\fi% select the smaller of both
  \ifdim\@tempa\p@<\p@\scalebox{\@tempa}{\usebox\pandoc@box}%
  \else\usebox{\pandoc@box}%
  \fi%
}
% Set default figure placement to htbp
\def\fps@figure{htbp}
\makeatother
\setlength{\emergencystretch}{3em} % prevent overfull lines
\providecommand{\tightlist}{%
  \setlength{\itemsep}{0pt}\setlength{\parskip}{0pt}}
\usepackage{bookmark}
\IfFileExists{xurl.sty}{\usepackage{xurl}}{} % add URL line breaks if available
\urlstyle{same}
% fallback for those not using the hyperref driver hyperxmp:
\makeatletter
\@ifundefined{xmpquote}{\newcommand{\xmpquote}[1]{#1}}{}
\makeatother
\hypersetup{
  hidelinks,
  pdfcreator={LaTeX via pandoc}}

\author{}
\date{}

\begin{document}

{
\setcounter{tocdepth}{2}
\tableofcontents
}
\section{WebCoder --- User Guide
(learners)}\label{webcoder--user-guide-learners}

\begin{quote}
One-paragraph summary: WebCoder is a bilingual (Romanian/English, plus
10 more interface languages) competitive-programming platform: browse
\textasciitilde1 800 problems mapped on the national curriculum, solve
them in a real code editor with instant judged verdicts, join rated
contests, climb the rating ladder, learn from hints and editorials, and
--- with a teacher --- do all of it as classroom homework. This guide
walks every learner flow end to end.
\end{quote}

\textbf{Who it\textquotesingle s for:} students, contest participants,
self-learners. Romanian twin:
\href{../guide-ro/GHID-UTILIZATOR.md}{\texttt{docs/guide-ro/GHID-UTILIZATOR.md}}.

\begin{center}\rule{0.5\linewidth}{0.5pt}\end{center}

\subsection{1. Create an account}\label{1-create-an-account}

\begin{enumerate}
\def\labelenumi{\arabic{enumi}.}
\tightlist
\item
  Open the site and press \textbf{Sign up}.
\item
  Choose a username and password (email optional unless your teacher
  requires verification; Google / GitHub / Microsoft school login also
  work).
\item
  That\textquotesingle s it --- you land on the home page with the
  catalog, contests and learning hub one click away.
\end{enumerate}

Under-16 accounts (GDPR Art. 8): profiles stay non-public and
school/county fields hidden until parental consent is recorded;
classroom accounts can be school-mediated instead.

\subsection{2. Find a problem}\label{2-find-a-problem}

The \textbf{Problems} catalog (\texttt{/problems}) is the main surface:

\begin{itemize}
\tightlist
\item
  filter by difficulty (very easy → hard), tag, curriculum chapter
  (grades 5--12 tree), source archive, I/O mode, solved-status;
\item
  the \textbf{Ctrl-K palette} searches problems, threads, articles,
  contests and users from anywhere;
\item
  the dice button picks a random problem matching your filters;
\item
  the catalog shows solve counters and your solved/attempted marks.
\end{itemize}

\pandocbounded{\includegraphics[keepaspectratio,alt={Catalog}]{C:/Users/Prekzursil/.zcode/workspace/draft-9cac9da0-09ed-462a-8bd2-a78591461ba4/wt-polish/docs/dist/USER-GUIDE/en/59810f9277299b0c9296bb4be89dab11b64bfb06.png}}

Curriculum view: pick your grade → chapter → topic; every node lists its
problems (pbinfo-style navigation, bilingual).

\subsection{3. Read, solve, submit}\label{3-read-solve-submit}

A problem page has the statement (Markdown + LaTeX math), limits,
samples with copy buttons and localized sample comments, hints, and the
editor.

\begin{enumerate}
\def\labelenumi{\arabic{enumi}.}
\tightlist
\item
  Pick a language --- C, C++, Python 3, Java, JavaScript, Go, C\#,
  Pascal (per-problem allowlists may apply).
\item
  Write code in the CodeMirror editor (Tab indents, Ctrl+Enter submits;
  your draft autosaves per problem and follows you across devices).
\item
  Optional before submitting: \textbf{Custom test} (run on your own
  input), \textbf{Run samples} (compile once, judge every sample with
  mismatch hints), the pre-submit \textbf{lint} (bilingual pitfall
  warnings), and the \textbf{stress runner} for authors.
\item
  Press \textbf{Submit} and confirm --- the verdict page polls the judge
  and shows the per-test table, time/memory meters against the actual
  bars, and the "how close am I?" band.
\end{enumerate}

\pandocbounded{\includegraphics[keepaspectratio,alt={Problem
page}]{C:/Users/Prekzursil/.zcode/workspace/draft-9cac9da0-09ed-462a-8bd2-a78591461ba4/wt-polish/docs/dist/USER-GUIDE/en/e1f7428dd0b4eabb1ad16241b72600ec8309c5f0.png}}
\pandocbounded{\includegraphics[keepaspectratio,alt={Verdict}]{C:/Users/Prekzursil/.zcode/workspace/draft-9cac9da0-09ed-462a-8bd2-a78591461ba4/wt-polish/docs/dist/USER-GUIDE/en/a81e1e860356660d999488ccdece170c87a65f71.png}}

Verdicts: \textbf{AC} (accepted), \textbf{WA} (wrong answer --- sample
tests show an expected-vs-actual diff), \textbf{TLE/MLE} (time/memory,
with the limit that was applied), \textbf{RE} (runtime error,
signal-mapped), \textbf{CE} (compile error, with plain-language hints).
Partial scoring is per test or per subtask, depending on the problem.

After an AC: the editorial unlocks (if gated), the solutions pool lets
you compare your code with the fastest/shortest accepted solutions, and
you can share a solve card.

\subsection{4. Learn, don\textquotesingle t just
score}\label{4-learn-dont-just-score}

\begin{itemize}
\tightlist
\item
  \textbf{Hints} --- progressive, per problem; the weekly digest emails
  your attempted-but-unsolved problems with their first hint.
\item
  \textbf{Editorials} --- unlock on your own AC by default (peeking
  marks the solve as "with editorial").
\item
  \textbf{/learning hub} --- spaced-repeat review queue over your
  misses, skill graph, theory cards, mistake-pattern profile, a
  generated weekly practice plan.
\item
  \textbf{/drills} --- a daily warm-up set (read-the-code,
  fill-the-line, complexity drills).
\item
  \textbf{/visualizers} --- step-frame animations of sorting and binary
  search.
\item
  \textbf{/courses} --- structured lessons (e.g. the Scratch→Python
  bridge).
\item
  \textbf{/start} --- the guided first-solve ladder + onboarding
  checklist.
\end{itemize}

\subsection{5. Contests and ratings}\label{5-contests-and-ratings}

\begin{itemize}
\tightlist
\item
  \texttt{/contests} lists upcoming (with countdown + ICS calendar
  subscription), running, and past contests; register in one click.
\item
  Formats include IOI-style (sum of best per problem), ICPC penalty,
  windowed personal-timer contests for classrooms, team contests, school
  cups and long challenges.
\item
  During a contest: live scoreboard (frozen near the end), clarification
  queue to the jury, self-prediction of your final rank.
\item
  After: Glicko rating updates with color bands and divisions
  (promotions/demotion events annotated on your rating graph), medals,
  replay of any participant\textquotesingle s timeline, and upsolving on
  the archive with full feedback.
\item
  Integrity: AI-assistance declarations, multi-account detection,
  plagiarism scans --- rated contests assume unaided work (see the RO/EN
  contest rules page).
\end{itemize}

\subsection{6. Community}\label{6-community}

\begin{itemize}
\tightlist
\item
  \textbf{Forum} with per-problem discussion threads (spoiler-flagged
  inside problem threads), wiki, blogs, articles hub.
\item
  \textbf{Messaging} --- private 1:1 threads with reporting, block and
  rate protections.
\item
  \textbf{Community hub} (\texttt{/amas}): reputation tiers, awards,
  medalist AMAs, problem-request board.
\item
  Report buttons everywhere feed the moderation queue (SLA-tracked).
\end{itemize}

\subsection{7. Profile \& settings}\label{7-profile--settings}

\texttt{/profile} shows stats, the 90-day heatmap, streaks
(weekly-target framing), per-tag skill bars, rating graphs with
annotations, badges, certificates (atestats with public verification
UUIDs) and your mistake-pattern profile. Settings: display name/bio,
county/school (segmented leaderboards), timezone (contest times render
in your zone), notification digest frequency, a11y display modes
(dyslexia-friendly, high-contrast), dark/light theme.

\subsection{8. Mobile \& offline}\label{8-mobile--offline}

The PWA installs on phones; offline statement packs
(\texttt{Save\ offline} on a problem) keep the statement, limits,
samples and hints available without a network; drafts sync when
you\textquotesingle re back online.

\subsection{9. Keyboard shortcuts}\label{9-keyboard-shortcuts}

\texttt{Ctrl-K} command palette · \texttt{?} keyboard map overlay ·
\texttt{Ctrl+Enter} submit · the map lists everything else.

\subsection{10. When something breaks}\label{10-when-something-breaks}

\begin{itemize}
\tightlist
\item
  Judge queue slow? \texttt{/status} shows live queue depth and platform
  health.
\item
  Account issues: password reset by email; contact a moderator via any
  report.
\item
  Bug? The forum\textquotesingle s support category or the
  problem-request board is the place.
\end{itemize}

\end{document}
